% 第6章 仿真实验与分析

\chapter{仿真实验与分析}\label{ch:experiments}

\section{实验设置}\label{sec:exp_setup}

% 环境：PandaPickCube (MuJoCo)
% 机器人：7自由度Franka Panda，OSC控制10Hz，物理仿真500Hz
% 任务：抓取方块并抬升20cm
% 评估：每个偏差值50个episode，报告成功率
% 训练：SAC, batch=256, utd=2, 100K在线步, 50个演示episode

\section{编码器偏差退化验证（H1）}\label{sec:exp_h1}

% 假设：编码器偏差使无偏差策略性能退化
%
% 实验：训练无偏差策略（18D），在bias 0.0到0.3 rad下评估
%
% 结果表：
% | 偏差 (rad) | 成功率 | 平均步数 |
% |-----------|--------|---------|
% | 0.00      | 100%   | 13.5    |
% | 0.05      | 100%   | 13.5    |
% | 0.10      | 80%    | 55.9    |
% | 0.15      | 62%    | 62.5    |
% | 0.20      | 51%    | 99.8    |
% | 0.25      | 5%     | 184.5   |
% | 0.30      | 0%     | 191.1   |
%
% 图：退化曲线（成功率 vs 偏差）
% 分析：>0.1 rad引起严重退化，>0.25接近完全失败

\section{固定偏差策略（H2）}\label{sec:exp_h2}

% 假设：在单一固定偏差下训练的策略会过拟合
%
% 实验：在固定b=0.2 rad下训练（24D），全范围评估
%
% 结果：训练偏差下100%，但零偏差时降至83%（过补偿）
%
% POMDP解释：策略学到了b=0.2的条件策略，
% 不具备一般化的b自适应能力

\section{随机偏差策略（H3）}\label{sec:exp_h3}

% 核心实验：在 $b \sim U[0, 0.25]$ rad下训练（24D）
%
% 三种策略对比表：
% | 偏差  | 无偏差策略 | 固定偏差策略 | 随机偏差策略 |
% |------|----------|------------|------------|
% | 0.00 | 100%     | 83%        | 92%        |
% | 0.05 | 100%     | 87%        | 100%       |
% | 0.10 | 80%      | 99%        | 99%        |
% | 0.15 | 62%      | 100%       | 96%        |
% | 0.20 | 51%      | 100%       | 97%        |
% | 0.25 | 5%       | 100%       | 96%        |
% | 0.30 | 0%       | 99%        | 99%        |
%
% 图：三策略对比图
%
% 关键发现：
% 1. 均衡鲁棒性：全范围92-100%
% 2. 超范围泛化：0.30 rad（训练上限=0.25）仍99%
% 3. 无过补偿：零偏差92%（固定偏差策略仅83%）
% 4. 统一策略处理任意偏差

\section{观测空间消融}\label{sec:exp_ablation}

% 验证第3.4节可观测性分析
%
% | 观测  | 维度 | 缺少的关键信号 | 结果     | 理论预测     |
% |------|-----|---------------|---------|-------------|
% | 18D  | 18  | real_tcp+block | 42%均值  | 不可辨识     |
% | 21D  | 21  | real_tcp      | 训练失败  | 知目标不知己 |
% | 24D  | 24  | （完整）       | 92-100% | 近似可辨识   |
%
% 理论（第3.4节）与实验的直接对应：
% - 18D失败：单步观测无法区分不同b
% - 21D失败：知道目标位置但不知道真实自身位置
% - 24D成功：$\Delta\vect{x}$ = real_tcp - biased_tcp 提供了b的信号

\section{备份策略实验}\label{sec:exp_backup}

% S1（单障碍物）结果：
% - 200步存活率
% - 训练曲线
%
% S2（移动+静止）结果：
% - 与S1的对比
%
% 奖励V1 vs V2 vs V3：
% - 训练曲线对比，展示V1/V2的失败模式
% - V3的收敛过程
%
% DR消融：
% - 有DR vs 无DR
% - 不同噪声水平下的性能
%
% 运动模式鲁棒性：
% - 按模式评估（LINEAR, ARC, STOP_GO, PASSING）

\section{联合系统评估}\label{sec:exp_integrated}

% 任务策略 + 备份策略切换：
% - 场景：抓取任务过程中周期性人手入侵
% - 指标：任务成功率、碰撞率、恢复率
% - 展示无缝策略切换
